\lecture{11}{Mon 30 Oct 2023 10:43}{}

Last time, we proved \textbf{Blunenthal 0-1 Law}: $A \in \mathcal{F}_0^+ \implies P(A) \in \{0,1\} $.

We have
\[
	\limsup_{h \to 0+} \frac{|B_{t + h} - B_t|}{\sqrt{2 h \log h \frac{1}{h}} } \in \mathcal{F}_0^+
.\] 

Let $\tau = \inf \{t >0 : B_t > 0\} $, $\sigma = \inf \{t = 0: B_t = 0\} $.

\begin{theorem}
	$P(\tau = 0) = P(\sigma = 0) = 1$
\end{theorem}
\begin{proof}
	\[
		\{\tau = 0\}  = \bigcap_{n = 0} ^{\infty } 
		\{\exists \varepsilon>0: 0 < \varepsilon < \frac{1}{n}, B_\varepsilon > 0\} \in \mathcal{F}_0^+
	.\] 
	\[
		\frac{1}{2} = P(B_t > 0) \le  P(\tau < t)
	.\]
	(?) Let $t \to  0$, we have
	\[
		P(\tau =0) \ge  \frac{1}{2} \implies P(\tau = 0) = 1
	.\] 
	By intermediate value theorem, $P(\sigma = 0) = 1$. (?)
\end{proof}

Recall that if $B_t$ is standard BM, then $t B_{\frac{1}{t}}$ is also.

\begin{theorem}[Kolmogorov 0-1 Law]
	Let $\mathcal{G}_t = \sigma\left( B_s: s \ge t \right) $.
	Let $\mathcal{T}_B = \bigcap_{t\ge 0} \mathcal{G}_t$.
	Then, $A \in \mathcal{T}_B \implies P(A) \in \{0,1\} $; $\mathcal{T}_B = \mathcal{F}_0^+(\omega)$.
\end{theorem}

\subsubsection{Reflection Principle of BM}

Let $M_t = \sup_{0\le s\le t} B_s$ be the running maximum. 

\begin{theorem}[Reflection Principle]
	For all $a > 0$,
	\[
		P(M_t > a) = 2 P(B_t > a) 
	.\] 
\end{theorem}
We actually have
\[
		P(M_t > a) = 2 P(B_t > a) = P(\tau_a< t)= P(|B_t| > a)
	.\]
\begin{remark}
	
	If $\tau_a = \inf \{t: B_t = a\} $, then
	\[
		P(\tau_a < t) = P(M_t > a)
	.\] 
	Moreover, the sets are equal:
	\[
		\{\tau_a < t\}  = \{M_t > a\} 
	.\] 


From the reflection principle,
\begin{align*}
	P(\tau_a < t) &= 2 P(B_t > a)\\
	&= 2 P(\sqrt{t} B_1 > a) \\
	&= 2 P\left( B_a > \frac{a}{\sqrt{t} } \right)  \\
	&= 2\left( 1 - P\left( B_a < \frac{a}{\sqrt{t} } \right)  \right)  \\
	&= 2\left( 1 - \Phi\left( \frac{a}{\sqrt{t} } \right)  \right)  
.\end{align*}
Therefore
\[
	P\left( \tau_a \in dt  \right) 
	= 2 \Phi'\left( \frac{a}{\sqrt{t} } \right) a t^{- 3 / 2}
	= \frac{a}{\sqrt{2 \pi t^3} } e^{- \frac{a^2}{2t}}
.\] 
Which is a \textit{stable distribution.}


\end{remark}

\begin{remark}
	Recall, using exponential martingale, gives
	\[
		E\left( e^{- \lambda \tau_a} \right) 
		= e^{- a \sqrt{2 \lambda} }
	.\] 
\end{remark}

\begin{intuition}
	Why is this called ``reflection principle''? 
	Let $T$ be finite a.s. stopping time. Let $B_t$ be standard BM. Then 
	\[
		B_t^r := \begin{cases}
			B_t & t \le T \\
			2 B_T - B_j & t > T
		\end{cases}
	.\] 
	is standard BM.
	
	This is proved using martingales.
	Same argument as the proof that $B_{T + t} - B_T$ is BM.
	
\end{intuition}

\begin{proof}[Proof of Reflection Principle]
	Observe that $W_t = B_{T + t} - B_T$ is BM. Therefore, so is $- (W_t) = B_T - B_{T + t}$.
	
	We have a ``silly representation'' of BM:
	\[
		B_t = \begin{cases}
			B_t & t \le T\\
			B_T + W_{t - T} & t > T
		\end{cases}
	.\] 
	Now, what if we put $- W_t$ instead of $W_t$?
	\[
		B_t^r = \begin{cases}
			B_t & t \le T\\
			B_T - W_{t - T} & t > T
		\end{cases}
	.\]
	$B_t^r$ follows $B_t$ until $T$, and then reflects about  $B_T$.

	Now begin the proof:
	\[
		\{M_t > a\}  = \{M_t > a, B_t > a\}  + \{M_t > a, B_t < a\}  
	.\] 
	If $B_t > a$ then $M_t > a$.
	\[
		P(M_t > a) = P(B_t > a) + P(M_t > a, B_t < a)
	.\] 
	But $\{M_t > a, B_t < a\} = \{M_t > a, B_t^r > a\} $. So
	\[
		P(M_t > a) = P(B_t > a) + P(B_t^r > a)
		= 2(B_t > a)
	.\] 
\end{proof}
